\section{The induction and global scales}\label{sec:scales}

The induction carries three estimates for each pair of deterministic words.  Two auxiliary laws $\mu,\nu$ can be postprocessed by one kernel to approximate the genuine trace laws, and the likelihood score $d\nu/d\mu-1$ has a small moment.  These two estimates control distinguishability at the current length.  The third estimate is common mass between traces of two consecutive powers of the first word; it supplies the hidden orientation bits used to construct the next pair.

Here is the quantitative idea behind the scale choices.  Ignoring the small exponent margins, a block of depth $b$ with target length $N$ starts from atoms of length about $N^{2^{-b}}$ whose squared score is at most $N^{-2^{-b}d}$.  The low-degree analysis will either suppress a term by signed smoothing or force it to contain $2^b$ old-score factors.  Those factors together retain approximately the decay $N^{-d}$.  The joint locations of their sibling groups supply the additional decay $N^{-(b-1)/4}$.  Sections~\ref{sec:block}--\ref{sec:coefficients} prove this estimate with the margins used below.  They also verify the input strength needed to control high degrees and preserve consecutive-power overlap.

Fix $0<c<c_+<1/(4\log2)$, then choose a constant integer $B\ge2$ and $0<\zeta<1$ such that
\[
 \frac{(1-\zeta)((B-1)/4-.05)}{B\log2}>c_+>c.
\]
Set
\[
\begin{split}
 D &=\left\lfloor (1-\zeta)\frac{\log H_{\rm del}}{B\log 2}\right\rfloor,
 \qquad j_0=5\cdot 2^B,\qquad
 \epsilon=2^{-4B-100}/(D+1),\\
 b_i &=\begin{cases}2,&1\le i\le j_0,\\ B,&j_0<i\le D.\end{cases}
\end{split}
\]
We henceforth take $D>j_0$, as holds eventually.  For either allowed block depth $b$, define the contraction exponent
\begin{equation}\label{eq:contraction-exponent}
 a_0(b)=(1+\epsilon)2^{-b}-\epsilon.
\end{equation}
It is positive for all sufficiently large instances.  Define the target scales backwards by
\[
 N_D=n,\qquad N_{i-1}=N_i^{a_0(b_i)}\quad(1\le i\le D),
\]
and the score exponents by
\[
 d_0=2.20,\qquad
 d_i=d_{i-1}+(b_i-1)/4-.05\quad(1\le i\le D).
\]
Each preliminary depth-two block adds $.20$, so
$d_{j_0}=2.20+2^B$.  This fixed preliminary phase raises the seed exponent to the strength required by depth-$B$ blocks.  The remaining blocks produce the asymptotic gain in the theorem.  The nested-overlap proof and the high-degree score estimate will show separately where this stronger input is used.

The remaining parameters control the accumulated projection error and rare events.  Put
\[
 \Delta=n^{-100(d_D+B+1)},\qquad
 r_i=2\left\lceil 10^8 2^{10B}(d_D+B+1)
       \frac{\log n}{\log N_i}\right\rceil\quad(0\le i\le D).
\]
The even moment orders $r_i$ are nonincreasing with $i$, so a moment bound available at one level is still available at the next.  Their dependence on $\log n/\log N_i$ makes the exceptional probabilities small relative to the final-length error budget even at the smallest scale.

\begin{lemma}[Uniform scale margins]\label{lem:scale-margins}
For the parameters above,
\begin{equation}\label{eq:smallest-scale}
 \log N_0\gtrsim_B (\log n)2^{-BD}
 \gtrsim_B (\log n)H_{\rm del}^{-(1-\zeta)}
 =\Lambda H_{\rm del}^{\zeta}.
\end{equation}
Moreover,
\[
 \epsilon\log N_0\gtrsim_B
 \frac{(\log n)^\zeta}{1+\log\log n}
 \gg\log\log n+\log(d_D+B+1).
\]
For every fixed nonnegative integer $h$, uniformly over $0\le i\le D$,
\begin{equation}\label{eq:uniform-absorption}
 (D+1)^h\Lambda=o(\epsilon\log N_0),\qquad
 \exp\bigl(O_B((D+1)^h\Lambda)\bigr)=N_i^{o(\epsilon)}.
\end{equation}
\end{lemma}

\begin{proof}
For every allowed $b$,
$a_0(b)\ge2^{-B}(1-(2^B-1)\epsilon)$.  Multiplying these bounds over $D$ blocks loses only a factor bounded below by a positive constant depending on $B$, because $D\epsilon\le2^{-4B-100}$.  The definition of $D$ then proves \eqref{eq:smallest-scale}.  Since $H_{\rm del}\le\log n$ and $D+1=O_B(1+\log\log n)$, the same bound proves the second display.  Finally,
\[
 \frac{(D+1)^h\Lambda}{\epsilon\log N_0}
 \lesssim_B\frac{(D+1)^{h+1}}{H_{\rm del}^{\zeta}}
 \longrightarrow0,
\]
because $D=O_B(\log H_{\rm del})$.  Also $N_i\ge N_0$, proving the uniform exponential formulation.
\end{proof}

In particular, $\log r_i=O_B(D+1)$ uniformly in $i$.  Equation~\eqref{eq:uniform-absorption} will control factors involving the channel, the number of blocks, and the fixed-depth calculations.  Constants here and below may depend on $B,\zeta,c_+$.  Unless additional dependence is explicit, constants in $O$, $\asymp$, $\lesssim$, and similar notation are absolute or depend only on fixed numerical parameters, such as a fixed degree cutoff; they do not depend on $\delta,n,D,r_i$.  All estimates used at multiple levels are uniform over those levels as $n\to\infty$.  In particular, the working rate $\delta$ always has the value fixed from the original instance, not a value recomputed from $N_i$.

Distinct equal-length words each have an exclusive full-length trace of positive probability, so neither genuine trace law dominates the other.  This motivates placing the likelihood score on auxiliary laws and recording their projection errors separately.

\begin{definition}[The level invariant]\label{def:invariant}
Choose $C>2$ and $0<\omega<1/8$ depending only on $\delta$, with
\[
 \log C=O(\Lambda),\qquad \log(\omega^{-1})=O(\Lambda),
\]
and a constant $M_B\ge0$ depending only on $B$.  The same choices are used at every level.
For $0\le i\le D$, a level-$i$ pair consists of two different strings $x_i,y_i$ of the same length $\ell_i\in[N_i/2,N_i]$, and probability laws $\mu_i,\nu_i$ on a common auxiliary space, which may depend on $n$ and $\delta$, satisfying the following conditions.
\begin{itemize}
\item There are a Markov kernel $\mathcal K_i$ from the auxiliary space to finite binary words and a number $E_i\ge0$ such that
\[
 \TV(\mathcal K_i\mu_i,\Tr(x_i))\le E_i,\qquad
 \TV(\mathcal K_i\nu_i,\Tr(y_i))\le E_i,
 \qquad \frac{E_i}{\ell_i}\le(i+1)\Delta.
\]
\item We have $\nu_i\ll\mu_i$ and
\[
 \left\|\frac{d\nu_i}{d\mu_i}-1\right\|_{L^{r_i}(\mu_i)}^2
       \le N_i^{-d_i}.
\]
\item For an integer $k_i\in[\ell_i,C^{2Bi+M_B}\ell_i]$,
\[
 \ov\bigl(\Tr(x_i^{k_i}),\Tr(x_i^{k_i+1})\bigr)\ge\omega.
\]
Here powers concatenate copies of the word.
\end{itemize}
\end{definition}
For every positive integer $m$, the first two conditions give
\[
 \TV\bigl(\Tr(x_i)^{\otimes m},\Tr(y_i)^{\otimes m}\bigr)
 \le\sqrt{mN_i^{-d_i}}+2mE_i
\]
by Hellinger tensorization and telescoping the projection errors, as detailed in Section~\ref{sec:transfer}.  This explains why the squared score carries the sample exponent $d_i$.

The pairs need not hit the target scales exactly.  Padding is reserved for the final transfer and is not inserted into words used by later blocks.

\begin{proposition}[Construction at every scale]\label{prop:construction}
Fix $B,\zeta,c_+$ and the parameters of this section.  Along every sequence with $n\to\infty$ and $\delta^3\log n\to\infty$, for all sufficiently large instances there exist level-$i$ pairs satisfying Definition~\ref{def:invariant} simultaneously for $0\le i\le D$.  The choices of $C,\omega$ depend only on $\delta$ and obey the displayed logarithmic bounds, and one $M_B$ depending only on $B$ works at every level.
\end{proposition}

The proof is by induction.  Section~\ref{sec:initial} supplies level zero.  Proposition~\ref{prop:block} constructs the next level, with its observation and likelihood estimates proved in Sections~\ref{sec:layout}--\ref{sec:coefficients}.  Section~\ref{sec:transfer} applies the final common kernel and returns to complete traces.
